\documentclass[10pt,a4paper]{amsart}
\usepackage[margin=19mm]{geometry}
\usepackage{amsmath,amssymb,mathtools}
\usepackage[scr=rsfs,cal=euler]{mathalfa}
\usepackage{xfrac}
\usepackage[unicode,hidelinks,bookmarks=false]{hyperref}
\newcommand{\ie}{i.e.\ }
\newcommand{\cf}{cf.\ }
\newcommand{\resp}{respectively\ }
\newcommand{\Subsection}{\subsection}
\providecommand{\nobreakdash}{\nobreak\hbox{-}}
\setlength{\emergencystretch}{2em}
\multlinegap0pt
\usepackage{bm}
\usepackage{bbm}
\usepackage{dsfont}
\usepackage{xspace}
\usepackage{cancel}
\usepackage{mathtools}
\usepackage{amsthm}
\makeatletter
\@ifundefined{theorem}{\newtheorem{theorem}{Theorem}}{}
\@ifundefined{definition}{\newtheorem{definition}[theorem]{Definition}}{}
\@ifundefined{lemma}{\newtheorem{lemma}[theorem]{Lemma}}{}
\makeatother
\title[Elliptic asymptotic behaviour of $q$-Painlev\'e transcendent]{Elliptic asymptotic behaviour of $q$-Painlev\'e transcendents}
\author{Joshua Holroyd}
\date{Source: arXiv:2506.05724v1 (CC BY 4.0)}
\begin{document}
\maketitle

\noindent\emph{Source and licence.} Elliptic asymptotic behaviour of $q$-Painlev\'e transcendents. Version arXiv:2506.05724v1 (CC BY 4.0), distributed under the Creative Commons Attribution 4.0 International licence, as recorded in the arXiv licence metadata for this exact version and shown on the abstract page https://arxiv.org/abs/2506.05724v1. Full rights notice: licenses/arxiv-2506.05724-CC-BY-4.0.txt.\par\smallskip

\noindent\textit{Editorial scope.} Counted unit: the Theorem whose source label is \texttt{thrm:main2} in the article (source0.tex lines 318--327, source label \texttt{thrm:main2}), delivered together with the article's own complete original proof of it (source0.tex lines 328--371, one \texttt{proof} environment of 44 lines). No mathematics is written, repaired or re-derived: statement and proof are the article's own text line for line. Required context retained uncounted: def:E (lines 90--96); eq:LOeq (lines 147--149); lem:cononForm (lines 157--170); thrm:elliptic (lines 178--188); def:cont (lines 209--215); int:Omega (lines 246--248); lem:diffEn0 (lines 283--292). Every definition, display and label of those lines is retained unchanged. Omitted from this package and not redistributed: 1--89, 97--146, 150--156, 171--177, 189--208, 216--245, 249--282, 293--317, 372--412. Nothing inside a retained excerpt is omitted or altered: every retained body is delivered exactly as the article's lines are written, so no omission receipt of any kind accompanies this package. Citations: the retained excerpts carry no citation command at all, so no bibliography entry of the article is transcribed and no reference is redistributed. Reference adaptation: none was needed and none was made; every reference inside the delivered unit resolves inside it, so no unresolved reference or citation is produced. Printed slips kept literal: no printed word, symbol, hypothesis or number of the counted statement or of its proof was altered. Layout and class: the article's own class is \texttt{amsart} with packages including biblatex, amsmath, amsthm, amssymb, float, enumitem, todonotes, color, graphicx; the wrapper is the lane's pinned \texttt{amsart} editorial wrapper at 10pt on A4, which restates the theorem-like environments the retained text needs and adds the article's own macro definitions that the retained text needs (none), with extra wrapper packages (bm, bbm, dsfont, xspace, cancel, mathtools). The delivered numbering is the unit's own. Assets: no retained span contains a figure environment, a graphics-inclusion command, a PDF-page inclusion, a table or a TikZ picture; no image file is copied into this package, the package declares no asset dependency and no image file is redistributed with it. Licence: version arXiv:2506.05724v1 of this article is distributed under the Creative Commons Attribution 4.0 International licence, as recorded in the arXiv licence metadata for this exact version and shown on the abstract page https://arxiv.org/abs/2506.05724v1. A full rights notice is delivered at licenses/arxiv-2506.05724-CC-BY-4.0.txt. Characters: every retained excerpt is pure ASCII, so the delivered engine, XeLaTeX with its default Latin Modern text font, reports no missing character.
\begin{definition}\label{def:E}
    We define the sequences $E:\mathbb{N}\rightarrow\mathbb{C}$ and $L:\mathbb{N}\rightarrow\mathbb{C}$ by
    \begin{align}\label{eq:invariant}
        E_n=&\frac{f_{n+1}^2f_n^2+t_0f_{n+1}f_n(f_{n+1}+f_n)+at_0(f_{n+1}+f_{n})+a}{f_{n+1}f_n},\\
        L_n=&\frac{(a+f_{n+1}f_n)(f_{n+1}+f_n)}{f_{n+1}f_n}.
    \end{align}    
\end{definition}

\begin{align}\label{eq:LOeq}
    f_{n+1}^2f_n^2+t_0f_{n+1}f_n(f_{n+1}+f_n)-E_0f_{n+1}f_n+at_0(f_{n+1}+f_{n})+a=H_n,
\end{align}

\begin{lemma}\label{lem:cononForm}
    With conditions $E_0\neq 2\pm4t_0$, define the sequence $F:\mathbb{N}\rightarrow\mathbb{C}$ by
    \begin{align*}
        f_n=\frac{c+F_n}{c-F_n}\quad\text{with $c\in\mathbb{C}^*$ such that}\quad c^4=\frac{2-E_0-4t_0}{2-E_0+4t_0}.
    \end{align*}
    Then considering bounded $n\in\mathbb{N}$ and taking the case $a=1$, Equation (\ref{eq:LOeq}) becomes
    \begin{align*}
        F_{n+1}^2F_{n}^2+\gamma\left(F_{n+1}^2+F_{n}^2\right)+\zeta F_{n+1}F_n+1=K_n\to0,\quad\epsilon\to0,
    \end{align*}
    where
    \begin{align*}
        \gamma=\frac{(2+E_0)c^2}{2-E_0-4t_0},\quad\zeta=\frac{8c^2}{2-E_0-4t_0},\quad K_n=\frac{(c-F_n)^2(c-F_{n+1})^2}{2-E_0-4t_0}H_n.
    \end{align*}
\end{lemma}

\begin{theorem}\label{thrm:elliptic}
    Consider bounded $n\in\mathbb{N}$ and  $a=1$. Regarding the sequence $F_n$ described in Lemma \ref{lem:cononForm} we obtain
    \begin{align*}
        F_n\to A\,\text{sn}(z_0+pn;k),\quad\epsilon\rightarrow0.
    \end{align*}
    Here, sn$(z;k)$ is Jacobi's elliptic sine function of argument $z$ and modulus $k$. Furthermore, $A^2=k$ and $k$ satisfying $0<|k|<1$ is uniquely determined by
    \begin{align*}
        k^2+\beta k+1=0,\quad \beta:=\frac{4+4\gamma^2-\zeta^2}{4\gamma}.
    \end{align*}
    Finally, $p$ is chosen such that $k\text{sn}^2(p;k)=-1/\gamma$ and $z_0$ is representative of the initial condition $f_0$, satisfying $A\text{sn}(z_0;k)=F_0$.
\end{theorem}

\begin{definition}\label{def:cont}
    Following Lemma \ref{lem:cononForm} and Theorem \ref{thrm:elliptic}, we define the continuous function $f:\mathbb{R}\rightarrow\mathbb{C}$ so that $f(n)=f_n$ for bounded $n\in\mathbb{N}$, and
    \begin{align*}
        f(x)\to\frac{c+A\,\text{sn}(z_0+px)}{c-A\,\text{sn}(z_0+px)},\quad\epsilon\rightarrow0,
    \end{align*}
    for bounded $x\in\mathbb{R}$. We likewise define corresponding continuous functions $E(x)$ and $L(x)$ (see Definition \ref{def:E}).
\end{definition}

        \begin{align}\label{int:Omega}
        \Omega(k):=\oint\frac{\text{d}u}{\sqrt{(1-u^2)(1-k^2u^2)}}.
    \end{align}

\begin{lemma}\label{lem:diffEn0}
    For any $n\in\mathbb{N}$ we have
    \begin{align}\label{eq:difference}
        E_n-E_{0}=&-\left(nL_n-\sum_{k=0}^{n-1}L_k\right)t_0\epsilon+M_n,
    \end{align}
    where
    \begin{align*}
        |M_n|<2|t_0|J_n\left(e^{|\epsilon|n}-1-|\epsilon|n\right)=\mathcal{O}\left(|\epsilon|^2n^2\right),\quad|\epsilon|n\ll1.
    \end{align*}
\end{lemma}

\begin{theorem}\label{thrm:main2}
    Let $a=1$ and assume $|L^{(k)}(0)|\leq R$ for all $k\in\mathbb{N}$ and some finite $R\in\mathbb{R}^+$. Furthermore, suppose we have $\eta\in\mathbb{N}$ such that $|\eta-\Omega/p|\ll1$, where $\Omega(k)$ is a period of the leading-order elliptic function given by Integral \eqref{int:Omega}. Then, regarding the difference $E_\eta-E_0$ we obtain
    \begin{align*}
        E_\eta-E_{0}\sim&-\frac{t_0\epsilon}{p}\left\{(4+L_0)\Omega(k)-8\Pi(k/c^2,k)\right\}+S_\eta,\quad\epsilon\to0,
    \end{align*}
    where $\Pi(\alpha^2,k)$ is the complete elliptic integral of the third kind with modulus $k$ and parameter $\alpha^2$, and
    \begin{align*}
        |S_\eta|<R|t_0\epsilon|\left\{|\eta-\Omega/p|e^{|\eta-\Omega/p|}+\left(\left|\Omega/p\right|+1\right)\left(e^{|\eta-\Omega/p|}-1\right)\right\}=\mathcal{O}\left(|\epsilon||\eta-\Omega/p|\right).
    \end{align*}
\end{theorem}

\begin{proof}
    For finite $\eta\in\mathbb{N}$, in Lemma \ref{lem:diffEn0} we showed that
    \begin{align*}
        E_\eta-E_{0}\sim&-\left((\eta+1)L_\eta-L_0-\sum_{k=1}^{\eta}L_k\right)t_0\epsilon,\quad\epsilon\to0.
    \end{align*}
    We apply the Euler-Maclaurin formula to the sum in the above equation, seeing that
    \begin{align*}
        \sum_{k=1}^{\eta}L(k)=\int_{0}^{\Omega/p}L(x)dx+\int_{\Omega/p}^{\eta}L(x)dx+\sum_{k=1}^{\infty}\frac{B_k}{k!}\left(L^{(k-1)}(\eta)-L^{(k-1)}(0)\right),
    \end{align*}
    where $B_k$, for $k\in\mathbb{N}$, are the Bernoulli numbers. Now, we utilise the expansion
    \begin{align*}
        L^{(k)}(\eta)=&\sum_{j=k}^{\infty}\frac{L^{(j)}(\Omega/p)}{(j-k)!}(\eta-\Omega/p)^{j-k},\quad k\in\mathbb{N},
    \end{align*}
    also noting that $L^{(j)}(\Omega/p)\to L^{(j)}(0)$ as $\epsilon\to0$ for all $j\in\mathbb{N}$. Applying these concepts, we obtain
    \begin{align}\label{eq:rawInt}
        &E_\eta-E_{0}\sim-t_0\epsilon\left(\frac{\Omega}{p}L_0-\int_{0}^{\Omega/p}L(x)dx\right)+S_\eta,
    \end{align}
    where
    \begin{align*}
        S_\eta=&-t_0\epsilon\sum_{k=1}^{\infty}\frac{\theta_k}{k!}(\eta-\Omega/p)^{k},
    \end{align*}
    and
    \begin{align*}
        \theta_k=(k-1)L^{(k-1)}(0)+\left(\Omega/p+1\right)L^{(k)}(0)-\sum_{j=1}^{\infty}\frac{B_j}{j!}L^{(j+k-1)}(0).
    \end{align*}
    Using the assumption that $|L^{(k)}(0)|\leq R$ for all $k\in\mathbb{N}$, and also noting
    \begin{align*}
        \sum_{j=1}^{\infty}|B_j/j!|<1,
    \end{align*}
    we may verify that
    \begin{align*}
        |S_\eta|<R|t_0\epsilon|\left\{|\eta-\Omega/p|e^{|\eta-\Omega/p|}+\left(\left|\Omega/p\right|+1\right)\left(e^{|\eta-\Omega/p|}-1\right)\right\}\ll|\epsilon|,
    \end{align*}
    for $|\eta-\Omega/p|\ll1$. Regarding the integral in Equation \eqref{eq:rawInt}, recall that in the $a=1$ case we have simply
    \begin{align*}
        L(x)=f(x)+\frac{1}{f(x)}+f(x+1)+\frac{1}{f(x+1)},
    \end{align*}
    and the continuous behaviour of $f(x)$ as $\epsilon\to0$ is stated in Definition \ref{def:cont}. Then, since the integral is across a period of the leading-order elliptic behaviour, we find that
    \begin{align*}
        \int_{0}^{\Omega/p}L(x)dx=&\frac{4}{p}\oint\left(\frac{1+ku^2/c^2}{1-ku^2/c^2}\right)\frac{\text{d}u}{\sqrt{(1-u^2)(1-k^2u^2)}}\\
        =&\frac{4}{p}\left(2\Pi(k/c^2,k)-\Omega(k)\right),
    \end{align*}
    where $\Pi(k/c^2,k)$ is a standard notation for the complete elliptic integral of the third kind, with parameter $k/c^2$ and modulus $k$. Applying this to Equation \eqref{eq:rawInt} completes the proof.
\end{proof}

\end{document}
